% chapter3banner.tex
% Standalone chapter-opening page: badge + title + photo + objectives + intro.
% This page is entirely separate from chapter3.tex's lesson content.

\khmerornatetitle{ជំពូកទី៣៖ អគ្គិសនី}{៣. អគ្គិសនី}
\khmerbannerphoto{img/lawson.jpg}

% --- CHAPTER LIST (2-column grid, edit freely) ---
\begin{center}
    \begin{tabularx}{\textwidth}{@{} >{\raggedright\arraybackslash}X >{\raggedright\arraybackslash}X @{}}
  \khmercircnum{1}~ ដែននិងកម្លាំងម៉ាញេទិក
  &
  \khmercircnum{2}~ អាំងឌុចស្យុងអេឡិចត្រូម៉ាញេទិក
  \\[3pt]
  \khmercircnum{3}~ អូតូអាំងឌុចស្យុង
  &
  \khmercircnum{4}~ សៀគ្វីចរន្តឆ្លាស់
  \\[3pt]
  \khmercircnum{5}~ រលកអេឡិចត្រូម៉ាញេទិក
    \end{tabularx}
\end{center}

\vspace{0.15cm}

% --- Intro paragraphs ---
\raggedright
\setlength{\parindent}{1.2em}
នៅជុំវិញខ្លួនយើង
យើងតែងតែសង្កេតឃើញនូវបរិក្ខារជាច្រើនដែលអាចប្រើប្រាស់បានលុះត្រាតែមានការចូលរួមពីចរន្តអគ្គិសនី។ កាលពីថ្នាក់ទី១១ យើងបានសិក្សាពីអគ្គិសនីចរន្តជាប់ម្តងហើយ ដែលលើកនេះ យើងនឹងសិក្សាពីបាតុភូតដែលនាំឲ្យមាននូវចរន្តអគ្គិសនីដែលជាហេតុនាំឲ្យយើងអាចប្រើប្រាស់វាបាន ដូចជាការសាកថ្មទូរស័ព្ទ កុំព្យូទ័រដែលជាឧបករណ៍អ្នកទាំងអស់គ្នាកំពុងប្រើប្រាស់ដើម្បីមើលសៀវភៅនេះ ព្រមទាំងជាឧបករណ៍ដែលកំពុងតែប្រើប្រាស់ប្រចាំថ្ងៃ។ មេរៀននៅក្នុងជំពួកនេះនឹងចូលរួមជាគ្រឹះដ៏សំខាន់សម្រាប់អ្នកដែលចង់ចាប់យកជំនាញខាងវិស្វករអគ្គិសនីនៅថ្នាក់ឧត្តមសិក្សា។

\clearpage